%%%%%%%%%%%%%%%%%%%%%%%%%%%%%%%%%%%%%%%%%%%%%%%%%%%%%%%%%%%%%%%%%%%%%%
% Lớp 4 | Chủ đề 1. Ôn tập và bổ sung | Số chẵn, số lẻ
% Dạng 4D1?3-3: Suy luận chẵn lẻ (lật cốc, bắt tay, chia đôi)
%%%%%%%%%%%%%%%%%%%%%%%%%%%%%%%%%%%%%%%%%%%%%%%%%%%%%%%%%%%%%%%%%%%%%%


% nguồn: 2026-11-18_lop4_chu-de-1-tu-luan (nhập 18/11/2026)
\begin{bt}%[4D1H3-3]%[ma:246512d2]
% nguon: Toán tư duy lớp 4 – Bài 3 SGK Toán 4 KNTT
	Mẹ mua $12$ quả táo và $15$ quả cam. Mẹ muốn xếp tất cả số quả vào hai đĩa sao cho số quả ở hai đĩa bằng nhau. Mẹ có làm được không? Vì sao?
	\loigiai{
		Tổng số quả: $12 + 15 = 27$ (quả).\\
		Muốn chia đều vào hai đĩa thì tổng số quả phải là số chẵn. $27$ là số lẻ.\\
		Vậy mẹ không làm được.
	}
\end{bt}

% nguồn: 2026-11-18_lop4_chu-de-1-tu-luan (nhập 18/11/2026)
\begin{bt}%[4D1V3-3]%[ma:2da2966b]
% nguon: Toán tư duy lớp 4 – Bài 3 SGK Toán 4 KNTT
	Trên bàn có $7$ cái cốc đều úp miệng xuống. Mỗi lần, bạn được lật đúng $2$ cốc bất kì. Hỏi sau một số lần lật, có thể làm cho cả $7$ cốc đều ngửa lên được không?
	\loigiai{
		Lúc đầu có $0$ cốc ngửa (số chẵn).\\
		Mỗi lần lật $2$ cốc thì số cốc ngửa hoặc tăng $2$, hoặc giảm $2$, hoặc giữ nguyên (lật $1$ cốc úp thành ngửa và $1$ cốc ngửa thành úp).\\
		Vậy số cốc ngửa luôn là số chẵn: $0, 2, 4, 6, \dots$\\
		Muốn cả $7$ cốc ngửa thì số cốc ngửa phải là $7$, là số lẻ. Vậy không thể làm được.
	}
\end{bt}

% nguồn: 2026-11-18_lop4_chu-de-1-tu-luan-bo-sung (nhập 18/11/2026)
\begin{bt}%[4D1N3-3]%[ma:26f4efb1]
% nguon: Toán tư duy lớp 4 – Bài 3 SGK Toán 4 KNTT
	Lớp 4A có $25$ học sinh. Cô giáo xếp mỗi bàn $2$ bạn. Có bàn nào chỉ có $1$ bạn ngồi không? Cần ít nhất bao nhiêu bàn?
	\loigiai{
		$25$ là số lẻ nên không chia hết thành từng cặp $2$ bạn: $25 : 2 = 12$ (dư $1$).\\
		Xếp được $12$ bàn mỗi bàn $2$ bạn, còn $1$ bạn ngồi một bàn riêng.\\
		Vậy có $1$ bàn chỉ có $1$ bạn và cần ít nhất $12 + 1 = 13$ (bàn).
	}
\end{bt}

% nguồn: 2026-11-18_lop4_chu-de-1-tu-luan-bo-sung (nhập 18/11/2026)
\begin{bt}%[4D1V3-3]%[ma:d7252a3a]
% nguon: Toán tư duy lớp 4 – Bài 3 SGK Toán 4 KNTT
	Có $9$ tấm thẻ ghi các số $1, 2, 3, 4, 5, 6, 7, 8, 9$. Có thể chia $9$ tấm thẻ thành hai nhóm sao cho tổng các số ở hai nhóm bằng nhau không? Vì sao?
	\loigiai{
		Tổng các số trên $9$ tấm thẻ: $1 + 2 + \dots + 9 = (1 + 9) \times 9 : 2 = 45$.\\
		Nếu hai nhóm có tổng bằng nhau thì tổng chung phải bằng hai lần tổng một nhóm, tức là số chẵn.\\
		$45$ là số lẻ, nên không thể chia được.
	}
\end{bt}
